% Exact registered equality-free Hilbert adequacy; editorial proof bridge.
\section{Adequacy of the registered first-order Hilbert presentation}
\label{sec:p02-hilbert-fol-adequacy}

The calculus here is literally Sheet~\ref{sheet:hilbert-fol}, with the
dependency annotations of Definition~\ref{def:p02-hilbert-fol-dependencies}:
primitive \(\neg,\to,\forall,\exists\), A1--A3, modus ponens, the three
quantifier schemas, premise-sensitive generalization, and the displayed
existential-elimination rule. Equality, falsity, conjunction and disjunction
are not extra primitive rules. This editorial bridge supplies its local
adequacy; it is not an unchanged single-source presentation of Mendelson,
Kleene or Enderton \cite{Mendelson1997,Kleene1952,Enderton2001}.
Kleene's own postulates for the predicate calculus, including his quantifier
postulates 9--12, and his deduction theorem with the condition that free
variables be held constant are presented, with a proof, in
\cref{sec:p02-w02-p02-025-kleene}.

The construction uses the current cited propositional basis theorem,
Theorem~\ref{thm:p02-calculi-derivability-equivalence}, only for A1--A3 and
modus ponens. First-order completeness comes from the independently proved
Henkin treatment and its equality-free specialization identified in
Annex~\ref{annex:p7.calculi.4}. Neither dependency assumes the Hilbert-FOL
adequacy now being established.

\begin{hbconvention}{def:p02-hfa-setting}{Literal language and dependency support}
Let \(\mathcal L\) be a first-order signature. A formula of the Hilbert
fragment is built from equality-free relation atoms using
\(\neg,\to,\forall,\exists\); \(r=t\) is not an atom of this fragment.
Write \(\Gamma\derives_{\mathsf H_{\mathsf{FOL}}}A\) only for the
registered dependency-annotated sequence. An axiom line has empty support;
a premise line has its singleton support; MP takes the union; Gen and
existential elimination retain their premise line's support while checking
their respective variable conditions. A support is a finite subset of
the permitted premise set, not every formula written elsewhere in the
sequence. All semantic consequences in this section range over nonempty
\(\mathcal L\)-structures and assignments, with premises true together
under the same assignment.
\end{hbconvention}

\begin{hblemma}{thm:p02-hfa-propositional-kernels}{proved=annex:p2.hfa.kernels}{Propositional kernels with exact supports}
Any finite A1--A3/MP derivation on the propositional sheet can have its
propositional variables uniformly replaced by Hilbert-fragment formulas,
giving a derivation using only the same A1--A3/MP rows on the first-order
sheet. A replaced premise retains its corresponding premise support, an
axiom has empty support, and MP takes unions. Consequently every
classical \(\neg,\to\) tautology supplies an empty-support kernel,
using the cited propositional basis result just named. No generalization,
existential elimination, equality rule or substitution on arbitrary
fixed premises is involved in this operation.
\end{hblemma}

The two kernels that the deduction theorem and the discharge interface below
use, the identity \(B\to B\) and the permutation
\((B\to(C\to G))\to(C\to(B\to G))\), do not need the cited basis theorem:
they have explicit empty-support A1/A2/MP derivations,
\cref{thm:p02-w02-p02-025-identity,thm:p02-w02-p02-025-permutation}.

\begin{hbtheorem}{thm:p02-hfa-soundness}{proved=annex:p2.hfa.soundness}{Soundness for the literal Hilbert-FOL rules}
For every line \((A,D)\) in a registered derivation, \(D\entails A\).
Thus \(\Gamma\derives_{\mathsf H_{\mathsf{FOL}}}A\) implies
\(\Gamma\entails A\), for arbitrary premise sets. The conclusion covers
open formulas and checks both premise-sensitive quantified rules.
\end{hbtheorem}

\begin{hblemma}{thm:p02-hfa-discharge-interface}{proved=annex:p2.hfa.discharge}{Dependency-preserving premise discharge}
If a registered sequence from \(\Gamma\cup\{B\}\) ends in
\((A,D)\), there is one from \(\Gamma\) deriving \(B\to A\)
with support contained in \(D\setminus\{B\}\). Conversely MP with
the added premise \(B\) recovers \(A\). The construction includes
Gen and existential elimination: it is not an unrestricted
generalization convention. This supplies the stronger support interface
used by the proof translations below and is compatible with the existing
deduction theorem, Theorem~\ref{thm:p02-hilbert-fol-deduction}.
\end{hblemma}

\begin{hbdefinition}{def:p02-hfa-encoding}{A closed falsum encoding for the full equality-free language}
Suppose the signature has an equality-free atomic formula \(Q\).
List its finitely many free variables in their fixed variable order and
let \(T\) be the universal closure of \(Q\to Q\). Thus \(T\) is a
closed Hilbert-fragment theorem; put \(F\defeq\neg T\).
Define the metalinguistic map \(h\) from full equality-free formulas
by preserving atoms, primitive negation, implication and both quantifiers,
and by
\[
 h(\bot)=F,\qquad
 h(B\land C)=\neg(h(B)\to\neg h(C)),\qquad
 h(B\lor C)=\neg h(B)\to h(C).
\]
Apply \(h\) memberwise to premise sets. This is a grammar translation,
not a formula identity or a new calculus. The closedness of \(F\)
prevents the encoding from introducing new free variables into
eigenvariable contexts. If the signature has no equality-free atom,
its literal Hilbert fragment contains no formula at all. For a
full-language translation in that case explicitly extend it by one
fresh nullary predicate and use that predicate for \(Q\); reduct and
expansion are part of this exceptional translation, not silent changes
to the literal sheet.
\end{hbdefinition}

\begin{hblemma}{thm:p02-hfa-alpha-encoding}{proved=annex:p2.hfa.alpha}{Bound renaming and the substitution interface}
For Hilbert-fragment formulas \(A=_{\alpha}B\), each direction
\(A\to B\) and \(B\to A\) has an empty-support registered derivation.
Thus alpha-conversion is derived, not a primitive sheet rule.
For the map \(h\),
\[
 \FV(h(A))=\FV(A),\qquad
 h(A[t/x])=_{\alpha}h(A)[t/x].
\]
It also preserves alpha-equivalence and classical satisfaction under
every structure and assignment. The relation \(=_{\alpha}\) is the
source's bound-name equivalence, not elementary equivalence or identity.
Every instance uses the actual total capture-avoiding substitution.
\end{hblemma}

\begin{hblemma}{thm:p02-hfa-existential-duality}{proved=annex:p2.hfa.exists-duality}{Primitive existential and its defined-negation counterpart}
For every fragment formula \(A\), each direction between
\(\exists xA\) and \(\neg\forall x\neg A\) has an empty-support
derivation on the literal registered sheet. This is a derived bridge,
not a definition of its primitive existential quantifier. The proof
uses only its quantifier schemas, the exact existential rule, Gen,
and the propositional kernels, with every variable restriction checked;
it does not invoke first-order completeness.
\end{hblemma}

\begin{hbtheorem}{thm:p02-hfa-proof-bridges}{proved=annex:p2.hfa.bridges}{Two explicit proof translations}
Every registered Hilbert derivation of \(A\) from \(\Gamma\)
translates to a derivation on the equality-free restriction of
Sheet~\ref{sheet:nd-fol}, with the same conclusion and open assumptions
contained in the final line's support. Conversely every finite
equality-free derivation on that ND sheet translates under \(h\)
to a registered Hilbert derivation of the translated conclusion from
the translated permitted assumptions. The construction checks every
propositional rule, both quantifier introduction and elimination pairs,
alpha-conversion, nonvacuous ND discharge, and the two Hilbert quantified
dependency restrictions. No equality inference is translated or
silently discarded.
\end{hbtheorem}

\begin{hbtheorem}{thm:p02-hfa-adequacy}{proved=annex:p2.hfa.adequacy}{Finite and strong adequacy of the literal registered fragment}
For any Hilbert-fragment \(\Gamma,A\),
\[
 \Gamma\derives_{\mathsf H_{\mathsf{FOL}}}A
 \quad\Longleftrightarrow\quad\Gamma\entails A.
\]
For arbitrary signatures and arbitrary premise sets this is proved in
classical ZFC, matching the inherited strong-completeness metatheory.
For a finite premise set it has a classical ZF proof, even in an
arbitrary signature: reduce the finite input to its finite occurring
signature and use the countable-language construction. The same ZF
proof applies to arbitrary premise sets in a supplied countable
signature. Every derivation has finite premise support. For the
full equality-free language with the stated encoding,
\[
 \Gamma\entails A
 \quad\Longleftrightarrow\quad
 h[\Gamma]\derives_{\mathsf H_{\mathsf{FOL}}}h(A),
\]
with the explicit fresh-predicate expansion when the original signature
has no equality-free atom. This is not adequacy for logical identity.
\end{hbtheorem}

The equality boundary is consequential. A signature may have function and
constant symbols but no relation symbols, so merely calling it
``first-order'' does not supply an atom for a bottom-free fragment.
Likewise adding equality atoms while retaining no identity schemas would
change the language and the completeness question. The equality-bearing
ND and LK treatments remain separate; their reflexivity and replacement
cannot be imported into this sheet by a reference to the new adequacy.
